% Imporditud: tools/import_eksperimendid.py
% Allikas: Eesti füüsikaolümpiaadi eksperimendi ülesannete
%   kogu (Rael Kalda, 2023), ülesanne Ü97, lahendus L97.
\setAuthor{EFO žürii}
\setRound{piirkonnavoor}
\setYear{2007}
\setNumber{G E2}
\setDifficulty{3}
\setTopic{staatika}
\setVahendid{täisnurkne kolmnurkne joonlaud nurkadega 30 ja 60 kraadi, kirjutuslaud}
\setPunktid{10}
\setAllikas{Eksperimendi kogu Ü97, lahendus lk 75}
\setLahendusseis{toores}

\prob{Kolmnurk}
Leida kolmnurga massikeskme koordinaadid, kui kolmnurga pikem kaatet on xtelg, lühem kaatet on y-telg ja nullpunktiks on joonlaua täisnurk.

\hint

\solu
1) Asetame kolmnurga laua servale, nii et skaalaga kaatet ristuks laua servaga. Nüüd liigutame kolmnurka üle laua serva, kuni kolmnurk kukub üle laua ääre. Siin tuleb leida viimane punkt, kus joonlaud veel ei kuku alla. Nüüd teame massikeskme x-koordinaati (eeldasime, et gradueeritud on pikem kaatet; kui skaala on lühemal kaatetil, siis saime teada massikeskme y-koordinaadi). (2 p)

A B

C

2) Kuna kolmnurga teine kaatet on skaalata, siis ei saa massikeskme teist koordinaati leida analoogiliselt x-koordinaadi määramisele. Nüüd asetame kolmnurga laua servale, nii et hüpotenuus on paralleelne laua servaga. Liigutame kolmnurka üle laua serva, kuni kolmnurk kukub üle laua ääre alla. Siin tuleb jällegi leida viimane punkt, kus joonlaud veel ei kuku alla. Kolmnurga skaalaga küljelt leiame, kui kaugel on massikese nullpunktist. (3 p) Leiame massikeskme y-koordinaadi selle järgi, kui kaugel asub massikese nullpunktist. Teeme joonise ja kanname sinna kaks joont, millest esimene kujutab massikeskme x-koordinaati ja teine kaugust hüpotenuusist. Esimese joone lõikepunkt x-teljega on punkt B ja teise joone lõikepunkt x-teljega on punkt A. Massikese asub nende kahe joone lõikepunktis (1 p). Näeme, et tekkis joonlauaga sarnane kolmnurk ABC, mille lühem külg BC ongi massikeskme y-koordinaat: BC = AB tan 30$^\circ \approx$ 0,57 AB. Lõigu AB pikkus on katses mõõdetud kauguste erinevus. (2 p)

Märkus: Abivahendeid kasutades saaks y-koordinaati määrata ka teisiti: joonistame kolmnurga skaala paberile. Asetame kolmnurga laua servale, nii et skaalaga kaatet oleks

paralleelne laua servaga. Toimides nii, nagu osas 1), ning kasutades paberit mõõtevahendina, leiame y-koordinaadi. Asetame kolmnurga laua servale, nii et skaalaga kaatet oleks paralleelne laua servaga. Liigutame kolmnurka üle laua serva ja leiame viimase punkti, kus joonlaud veel ei kuku alla. Märgime vastava koha laual sõrmega ja mõõdame kolmnurga abil sõrme kauguse laua servast. See annabki y-koordinaati. Selline sõrme kasutamine ei kuulu tema tavapäraste funktsioonide hulka, vaid ta esineb abivahendi rollis. Taoliste meetodite korral anda kogu osa 2) eest ainult 2 punkti.

3) Katsete kordamine vähemalt kolm korda --- 1 p. 4) Õige ja üheselt mõistetav vastus --- 1 p.
\probend
